% disclaimer_legal_privacy.tex --- Appendix A
\appendix
\section{Disclaimer, Legal, and Privacy Policy}
\label{sec:disclaimer}

This appendix documents the ethical, legal, and privacy framework under
which the research was conducted. These policies were in effect throughout
the project and are enforced structurally in the released artifacts.

\subsection{AI-assisted authorship disclaimer}
\label{sec:disclaimer:ai}

Parts of this manuscript, the accompanying code, and the blog posts at
\url{https://pocoo.vaked.dev} were drafted with assistance from AI models
(Claude, GLM, Qwen2.5, DeepSeek). The research direction, experimental
design, evaluation criteria, and all scientific claims were conceived and
verified by the human author. AI assistance was used for prose synthesis,
code generation, and literature verification, consistent with the Loop
Engineering paradigm in which the human designs the system and the system
executes \citep{osmani2026loop}. As noted in the project's public writing:
``the disclaimer is part of the loop too''---the AI that assisted is itself
a component of the system being described.

\subsection{Data privacy and PII scrubbing}
\label{sec:disclaimer:privacy}

The \texttt{ultrawhale-dogfood} dataset enforces PII scrubbing as a
structural invariant. The dataset schema includes a boolean field
\texttt{pii\_scrubbed} that is set to \texttt{true} for every published
record. The dogfeed loop that generates training data from LLM interactions
scrubs personally identifiable information (names, email addresses, phone
numbers, API keys, credentials) before depositing records to HuggingFace.
The dataset is labeled \texttt{Synthetic} on HuggingFace and contains no
human-generated personal data. All training pairs are either synthetic
(generated by the dogfeed loop from free-tier LLM responses) or
deterministically generated (the domain generators in
\texttt{build\_domain\_data.py} produce seeded, reproducible pairs with no
real-world data).

\subsection{Genesis contract and provenance signing}
\label{sec:disclaimer:genesis}

Every dataset release and model checkpoint carries a \emph{genesis
contract}: a structured declaration of what the artifact reduces, what makes
a valid iteration, when the loop stops, and what must never be sacrificed
\citep{vaked2024genesis}. The dataset schema includes \texttt{genesis} and
\texttt{genesis\_seal} fields that cryptographically bind each release to
its declared contract. The \texttt{produced\_by} and \texttt{signed} fields
provide end-to-end provenance: every record can be traced to the loop
iteration and model version that produced it. This makes the research
artifacts auditable---a property we consider necessary for open-source ML
research that operates outside traditional institutional review structures.

\paragraph{Repo-level genesis seal.} In addition to the dataset-level seal,
this repository carries a \emph{genesis seal} (\texttt{.genesis\_seal.json})
that cryptographically binds the current repo state to the genesis contract.
The seal is a SHA-256 hash over: (1) the genesis contract text in this
appendix, (2) the manuscript line count and table count, (3) the baseline
results JSON, and (4) the git HEAD commit hash. Any change to the manuscript,
baselines, or genesis contract invalidates the seal. The seal is verified
automatically in CI on every push and weekly via a scheduled DNS link check
(\texttt{.github/workflows/link-check.yml}). To verify locally:
\begin{verbatim}
python tools/genesis_seal.py --verify
\end{verbatim}
This is the honesty loop: the seal makes the research state auditable, the
CI makes the audit automatic, and the genesis contract makes the audit
meaningful by declaring what must never be sacrificed.

\subsection{Adversarial evaluation: safety note}
\label{sec:disclaimer:heretic}

The heretic adversarial evaluation (Section~\ref{sec:eval:paradox}) uses
prompts about dense technical content including chemical formulas, CVE
identifiers, exploit techniques, and pharmaceutical data. These prompts were
generated using the \texttt{heretic} tool \citep{heretic2024} and
Qwen2.5-7B prompt generation. No harmful content is produced or stored: the
prompts are used solely to generate \emph{responses} whose token distribution
is maximally dense in must-keep patterns (numbers, identifiers, formulas).
The evaluation measures token survival in compressed output, not the
semantic content of the responses. The compressed outputs are not published;
only aggregate metrics ($\mathit{exact\_keep\_pct}$, $\mathit{keep\_rate}$)
are reported. This is a compression-quality benchmark, not a content
generation system.

\subsection{Licensing}
\label{sec:disclaimer:license}

All software artifacts are released under Apache~2.0. The dataset
(\texttt{ultrawhale-dogfood}) is released under Apache~2.0 on HuggingFace.
The model checkpoints (\texttt{PeetPedro/kompress-v*}) are released under
Apache~2.0. The blog posts at \url{https://pocoo.vaked.dev} are public
technical writing. This manuscript is submitted to arXiv under the arXiv
license; the target venue (ICLR 2027) will determine the final publication
license.

\subsection{Decentralized funding and conflicts of interest}
\label{sec:disclaimer:funding}

This research was not supported by any institutional grant, corporate
sponsorship, or venture capital. The compute was funded by the author's
personal expenditure (\$37.19 DeepSeek API + \$0.55 vast.ai GPU) and a
decentralized open-source funding model in which a community council guides
donations (Section~\ref{sec:data}). The author declares no competing
financial interests. The open-source artifacts are released without
monetary restriction; the community council's role is limited to
infrastructure cost reimbursement and maintainer time support, with no
editorial control over research direction or claims.
